\documentclass[a4paper,12pt]{article}
% ============================================================
% Общий преамбул для отчётов по лабораторным работам (labs_lpr)
% Подключается из report/main.tex командой:
%   % ============================================================
% Общий преамбул для отчётов по лабораторным работам (labs_lpr)
% Подключается из report/main.tex командой:
%   % ============================================================
% Общий преамбул для отчётов по лабораторным работам (labs_lpr)
% Подключается из report/main.tex командой:
%   \input{../../common/latex_template/preamble.tex}
% ============================================================

% Поддержка русского языка и кодировки
\usepackage[T2A]{fontenc}
\usepackage{lmodern}
\usepackage[utf8]{inputenc}
\usepackage[russian]{babel}

% Математические пакеты
\usepackage{amsmath, amssymb}

% Настройка полей
\usepackage{geometry}
\geometry{top=2cm, bottom=2cm, left=2.5cm, right=1.5cm}

% Колонтитулы
\usepackage{fancyhdr}
\usepackage{graphicx}

% Современные таблицы
\usepackage{tabularray}

% Работа с csv файлами
\usepackage{csvsimple-l3}
\usepackage{tikz}

% Графики
\usepackage{pgfplots}
\pgfplotsset{compat=1.18}

% Ссылки
% \usepackage{hyperref} % гипперссылки, нужно оформление вместо красного квадрата
\usepackage[russian]{cleveref} % умные ссылки, нужно переоформление для формул, рисунков и тд (вместо ф-л что-то нужно придумать)

% ============================================================

% Поддержка русского языка и кодировки
\usepackage[T2A]{fontenc}
\usepackage{lmodern}
\usepackage[utf8]{inputenc}
\usepackage[russian]{babel}

% Математические пакеты
\usepackage{amsmath, amssymb}

% Настройка полей
\usepackage{geometry}
\geometry{top=2cm, bottom=2cm, left=2.5cm, right=1.5cm}

% Колонтитулы
\usepackage{fancyhdr}
\usepackage{graphicx}

% Современные таблицы
\usepackage{tabularray}

% Работа с csv файлами
\usepackage{csvsimple-l3}
\usepackage{tikz}

% Графики
\usepackage{pgfplots}
\pgfplotsset{compat=1.18}

% Ссылки
% \usepackage{hyperref} % гипперссылки, нужно оформление вместо красного квадрата
\usepackage[russian]{cleveref} % умные ссылки, нужно переоформление для формул, рисунков и тд (вместо ф-л что-то нужно придумать)

% ============================================================

% Поддержка русского языка и кодировки
\usepackage[T2A]{fontenc}
\usepackage{lmodern}
\usepackage[utf8]{inputenc}
\usepackage[russian]{babel}

% Математические пакеты
\usepackage{amsmath, amssymb}

% Настройка полей
\usepackage{geometry}
\geometry{top=2cm, bottom=2cm, left=2.5cm, right=1.5cm}

% Колонтитулы
\usepackage{fancyhdr}
\usepackage{graphicx}

% Современные таблицы
\usepackage{tabularray}

% Работа с csv файлами
\usepackage{csvsimple-l3}
\usepackage{tikz}

% Графики
\usepackage{pgfplots}
\pgfplotsset{compat=1.18}

% Ссылки
% \usepackage{hyperref} % гипперссылки, нужно оформление вместо красного квадрата
\usepackage[russian]{cleveref} % умные ссылки, нужно переоформление для формул, рисунков и тд (вместо ф-л что-то нужно придумать)

\newcommand{\qq}{\textbf{???}}
\InputIfFileExists{../data/processed/results.tex}{}{}
\newcommand{\res}[1]{\ifcsname res@#1\endcsname\csname res@#1\endcsname\else\qq\fi}
\newcommand{\mainmetal}{\ifcsname res@primary\endcsname\csname res@primary\endcsname\else Cu\fi}
\newcommand{\labimage}[2]{\IfFileExists{img/#1.pdf}{\includegraphics[width=#2\linewidth]{img/#1.pdf}}%
{\fbox{\parbox[c][2.4cm][c]{#2\linewidth}{\centering График появится после измерений}}}}
\pagestyle{fancy}\fancyhf{}
\fancyhead[L]{\includegraphics[height=1cm]{../../common/latex_template/miptlogo.png}}
\fancyhead[R]{Емельянов Д.\,А.}\fancyfoot[C]{\thepage}\setlength{\headheight}{33pt}
\begin{document}
\begin{titlepage}\centering
{\large Московский физико-технический институт\\
Физтех-школа физики и исследований им. Ландау}\par
\vfill
{\Large Лабораторная работа № 3.3.5\\ по курсу общей физики}\par\bigskip
{\huge\bfseries Эффект Холла в металлах}\par
\vfill
\begin{flushright}Емельянов Дмитрий Алексеевич\\Группа Б02-502\\
Соавтор: \qq\\Дата опыта: \qq\end{flushright}
\vfill г. Долгопрудный
\end{titlepage}

\section*{Аннотация}
По поперечному напряжению в металлической пластинке с током определяется постоянная Холла. Для основного образца (медь или серебро) используются серии при нескольких токах, для цинка --- одна серия при максимальном токе. Продольное падение напряжения даёт независимое измерение удельной проводимости. Совместно эти данные позволяют оценить холловскую концентрацию и подвижность носителей.

\section{Теоретические сведения}
Выберем ось $x$ вдоль тока образца $I$, ось $z$ вдоль прямого поля и ось $y$ так, чтобы система была правой. В стационарном состоянии поперечное электрическое поле уравновешивает магнитную составляющую силы Лоренца. Для одного типа носителей с зарядом $q$, концентрацией $n$ и дрейфовой скоростью $v_x$ имеем
\[
 q(E_y-v_xB_z)=0,\qquad j_x=nqv_x=\frac{I}{ah}.
\]
При ширине пластинки $a$ и толщине $h$ получаем
\begin{equation}
 U_H\equiv aE_y=R_H\frac{IB_z}{h},\qquad R_H=\frac{1}{nq}.
 \label{eq:hall}
\end{equation}
В этой системе $U_H$ равно потенциалу края при меньшем $y$ минус потенциал противоположного края. Для перехода от показаний прибора введём множитель $s=\pm1$, определяемый по положению контактов и полярности включения. Напряжения $U_{24}$ ниже приведены к одной полярности, в том числе после переключения К2.
\begin{equation}
 u=U_{24}(B)-U_{24}(0),\qquad U_H=s u,\qquad
 u_{\mathrm{odd}}(B)=\frac{U_{24}(+B)-U_{24}(-B)}{2}.
 \label{eq:odd}
\end{equation}
Вычитание нуля устраняет подмешивание продольного напряжения из-за смещения поперечных контактов. Полуразность при равных модулях поля дополнительно устраняет чётный по полю вклад. Продольное напряжение $U_{34}$ измеряется отдельно при $B=0$. Если $l$ --- расстояние между контактами 3 и 4, то
\begin{equation}
 \rho_0=\frac{|U_{34}|ah}{|I|l},\quad \sigma_0=\rho_0^{-1},\quad
 n_H=\frac{1}{e|R_H|},\quad \mu_H=|R_H|\sigma_0.
 \label{eq:transport}
\end{equation}
В многозонном металле $n_H$ и $\mu_H$ характеризуют холловский отклик. Отождествление их с микроскопическими параметрами носителей требует однозонного приближения.

\section{Установка и приборы}
Источник питания электромагнита задаёт ток $I_\text{М}$, амперметр А1 измеряет его, разъём К1 меняет направление поля. В независимой цепи образца ток задают источником и реостатом, отсчёт снимают по А2. Переключатель К3 выбирает контакт 2 или 3, контакт 4 общий. К2 меняет полярность входа микровольтметра. Поле измеряют миллитесламетром либо пробной катушкой с милливеберметром.
\begin{figure}[ht]\centering
\begin{tikzpicture}[>=stealth]
\draw[thick,fill=gray!10] (0,0) rectangle (6,1.5);
\draw[->,thick] (-1,.75)--(0,.75) node[midway,above]{$I$};
\draw[->,thick] (6,.75)--(7,.75);
\foreach \x/\y/\n in {2/1.5/2,2/0/4,4.8/0/3}
  {\fill (\x,\y) circle(2pt);\node at (\x,{\y+(\y>0 ? .23 : -.23)}) {\n};}
\node at (3.5,.8) {$\odot\, B_z$};
\draw[<->] (2,-.55)--(4.8,-.55) node[midway,below]{$l$};
\draw[<->] (-.42,0)--(-.42,1.5) node[midway,left]{$a$};
\end{tikzpicture}
\caption{Назначение контактов. Поле направлено к наблюдателю, толщина $h$ отсчитывается перпендикулярно рисунку. Для определения знака требуется зафиксировать ориентацию реального держателя.}\label{fig:contacts}
\end{figure}

\begin{table}[ht]\centering\small
\caption{Параметры, снимаемые с установки и паспортов приборов}
\begin{tblr}{colspec={l X[l]},width=.94\linewidth,hlines}
Величина & Запись в журнале\\
$h,a,l$ каждого образца, мм & Размеры с держателя и их погрешности: \qq\\
А1, А2 & Модель, диапазон, цена деления, класс: \qq\\
Микровольтметр & Пределы для $U_{24}$ и $U_{34}$, цена деления, класс: \qq\\
Измеритель поля & Тип, диапазон и погрешность: \qq; для катушки $SN$: \qq\ м$^2$\\
\end{tblr}
\end{table}

\section{Методика измерений}
\begin{enumerate}
\item Проверить регулировку тока медного или серебряного образца и диапазон тока магнита при К1 в положении I. Записать допустимые токи с установки.
\item При медленном увеличении $I_\text{М}$ снять не менее 15 точек во всём диапазоне. Измерять $B$ непосредственно либо по изменению потокосцепления при извлечении катушки: $B=|\Delta\Phi|/(SN)$. Значение $SN$ снять с держателя.
\item Вставить основной образец в зазор выключенного магнита. При фиксированном $I$ измерить $U_0=U_{24}(0)$. Снять 5--7 значений $U_{24}$ при плавном нарастании поля. Резкое изменение потока наводит ЭДС в измерительных проводах.
\item Повторить серию для 5--7 значений $I$, каждый раз заново измеряя $U_0$. При максимальном $I$ снять также обратную ветвь при тех же модулях поля. Для цинка измерить $U_0$ и одну серию при максимальном токе.
\item Для обоих держателей зарисовать ток, направление поля по стрелке на магните и положению К1, расположение контактов и полярность прибора. Отметить направление силы $q\mathbf v\times\mathbf B$ и край с повышенным потенциалом. Определить $s$ по \cref{fig:contacts} с учётом реального монтажа.
\item Вынуть каждый образец из зазора, перевести К3 на $U_{34}$ и при токе порядка максимального записать $I$ и $U_{34}$. При смене предела записать новую цену деления.
\end{enumerate}

\section{Результаты и обработка}
Поле для каждой точки находим линейной интерполяцией градуировки между соседними измерениями. Обратной ветви присваиваем $B_z<0$, предполагая равенство модулей полей при обращении тока. Для каждой серии подбираем прямую $u=kB_z+b$. Смещение $b$ служит проверкой вычитания нуля. Сравнение прямой и обратной ветвей при одинаковом $I$ проверяет воспроизводимость наклона.
\begin{figure}[ht]\centering
\labimage{calibration}{.70}
\caption{Градуировка электромагнита $B(I_\text{М})$ на возрастающей ветви.}
\end{figure}
\begin{figure}[ht]\centering
\labimage{hall_\mainmetal}{.72}
\caption{ЭДС Холла в полярности контактов основного образца при фиксированных токах. Обратная ветвь показана при отрицательных $B_z$.}
\end{figure}

\begin{table}[ht]\centering\small
\caption{Наклон $k$ и свободный член $b$ по сериям. Знак ветви задан положением К1.}
\ifcsname res@hasSlopes\endcsname
\csvreader[tabularray={colspec={ccccccc},hlines},table head={Образец & $I$, А & Ветвь & $k$ & $\delta k$ & $b$ & $\delta b$\\}]
{../data/processed/slopes.csv}{}{\csvexpval\csvcoli & \csvexpval\csvcolii & \csvexpval\csvcoliii & \csvexpval\csvcoliv & \csvexpval\csvcolv & \csvexpval\csvcolvi & \csvexpval\csvcolvii}
\else
\begin{tblr}{colspec={ccccccc},hlines}
Образец & $I$, А & Ветвь & $k$ & $\delta k$ & $b$ & $\delta b$\\
\qq & \qq & \qq & \qq & \qq & \qq & \qq
\end{tblr}\fi
\par\smallskip $k$ и $\delta k$ в мкВ/Тл, $b$ и $\delta b$ в мкВ.
\end{table}
Для основного образца по прямым ветвям определяем $K=dk/dI$ из линейной аппроксимации со свободным членом, для цинка используем $k/I$ одной серии. При $k$ в В/Тл, $I$ в А и $h$ в м из \cref{eq:hall} следует
\[
 R_H^{\mathrm{Cu/Ag}}=s h K,\qquad R_H^{\mathrm{Zn}}=s h k/I.
\]
\begin{figure}[ht]\centering
\labimage{slopes_\mainmetal}{.69}
\caption{Наклон $k(I)$ для основного образца. По второй ступени МНК $K=\res{\mainmetal Slope}$ мкВ/(Тл$\cdot$А).}
\end{figure}
\begin{figure}[ht]\centering
\labimage{hall_Zn}{.69}
\caption{Поперечное напряжение цинкового образца при максимальном токе.}
\end{figure}

\paragraph{Погрешности.}
При отсутствии паспортных данных для показания $X$ используем приборную составляющую $0{,}03|X|$ и отсчёт в одну единицу последнего записанного разряда. Для распространения принимаем их за стандартные неопределённости. Скрипт варьирует отсчёты независимо, а масштаб одного прибора, калибровочные точки, $U_0$ одной серии и размеры одного образца использует согласованно во всех зависимых величинах. В каждом из 4000 вариантов повторяется МНК; разброс результатов даёт абсолютную погрешность. Для интерполяции используется линеаризация на выбранном интервале. Такой расчёт описывает погрешности входных данных; возможный гистерезис магнита и нелинейность проверяют по остаткам и обратной ветви.

\clearpage
\section{Обсуждение}
\newcommand{\entry}[2]{\(\res{#1#2}\)}
\newcommand{\pct}[2]{\(\res{#1#2Pct}\)}
\begin{table}[ht]\centering\small
\caption{Параметры основного образца (\res{primary}) и цинка по \cref{eq:transport}. В столбцах $\delta$, \% приведена относительная погрешность.}
\begin{tblr}{colspec={l c c c c},hlines}
Величина & Основной образец & $\delta$, \% & Zn & $\delta$, \%\\
$R_H$, м$^3$/Кл & \entry{\mainmetal}{RH} & \pct{\mainmetal}{RH} & \entry{Zn}{RH} & \pct{Zn}{RH}\\
$n_H$, м$^{-3}$ & \entry{\mainmetal}{N} & \pct{\mainmetal}{N} & \entry{Zn}{N} & \pct{Zn}{N}\\
$\rho_0$, Ом$\cdot$м & \entry{\mainmetal}{Rho} & \pct{\mainmetal}{Rho} & \entry{Zn}{Rho} & \pct{Zn}{Rho}\\
$\sigma_0$, См/м & \entry{\mainmetal}{Sigma} & \pct{\mainmetal}{Sigma} & \entry{Zn}{Sigma} & \pct{Zn}{Sigma}\\
$\mu_H$, см$^2$/(В$\cdot$с) & \entry{\mainmetal}{Mu} & \pct{\mainmetal}{Mu} & \entry{Zn}{Mu} & \pct{Zn}{Mu}
\end{tblr}
\end{table}

Проверки по данным: совместимы ли $b$ с нулём и наклоны прямой и обратной ветвей между собой в пределах погрешностей: \qq.
В таблице 6 приложения практикума приведены следующие значения при $20\,^\circ$C:
\begin{center}
\begin{tblr}{colspec={lccc},hlines}
Металл & $R_H$, $10^{-10}$ м$^3$/Кл & $\sigma_0$, $10^7$ См/м & $\mu_H$, см$^2$/(В$\cdot$с)\\
Cu & $-0{,}55$ & 5,96 & 33\\
Ag & $-0{,}9$ & 6,28 & 57\\
Zn & $+0{,}55$ & 1,69 & 9,3
\end{tblr}
\end{center}
Для сравнения используем $\rho_{\rm табл}=1/\sigma_{\rm табл}$ и разность $X-X_{\rm табл}$. Температура опыта: \qq. Расхождения с учётом погрешностей: \qq.

В п.~19 методички поставлен вопрос об отрицательном $R_H$ цинка. В слабом поле для электронной и дырочной групп с концентрациями $n,p$ и подвижностями $\mu_e,\mu_h$ двухносительная модель даёт
\[
 R_H=\frac{p\mu_h^2-n\mu_e^2}{e(p\mu_h+n\mu_e)^2}.
\]
Отрицательный знак возможен при $n\mu_e^2>p\mu_h^2$, положительный --- при обратном неравенстве. Таблица 6 указывает для цинка положительный знак, что противоречит формулировке п.~19. Поэтому знак для данного держателя определяют по ориентации и измеренному наклону; одного подсчёта валентных электронов для его объяснения недостаточно.

\section{Заключение}
Для основного образца получено $R_H=\res{\mainmetal RH}$ м$^3$/Кл, для цинка $R_H=\res{ZnRH}$ м$^3$/Кл. Знак холловского отклика основного образца и цинка: \qq. Согласие транспортных параметров с таблицами и основное ограничение точности: \qq.

\smallskip\noindent\textit{Методика:} Лабораторный практикум по общей физике. Т.~2. Электричество и магнетизм. Работа 3.3.5, с.~183--187; введение к разделу III.
\end{document}
